\documentclass[11pt]{amsart}

\usepackage[margin=1in]{geometry}
\usepackage{amsmath,amssymb,mathtools}
\usepackage{array,booktabs}
\usepackage[T1]{fontenc}
\usepackage{lmodern}
\usepackage{microtype}
\usepackage{xcolor}
\usepackage{xurl}
\usepackage[colorlinks=true,
            linkcolor=blue!55!black,
            citecolor=blue!55!black,
            urlcolor=blue!55!black]{hyperref}

\newtheorem{theorem}{Theorem}[section]
\newtheorem{proposition}[theorem]{Proposition}
\newtheorem{lemma}[theorem]{Lemma}
\newtheorem{corollary}[theorem]{Corollary}
\theoremstyle{definition}
\newtheorem{definition}[theorem]{Definition}
\theoremstyle{remark}
\newtheorem{remark}[theorem]{Remark}

\newcommand{\Z}{\mathbf Z}
\newcommand{\F}{\mathbf F}
\newcommand{\Ga}{\mathbb G_a}
\newcommand{\Gm}{\mathbb G_m}
\newcommand{\Spec}{\operatorname{Spec}}
\newcommand{\Soc}{\operatorname{Soc}}
\newcommand{\gr}{\operatorname{gr}}
\newcommand{\length}{\operatorname{length}}
\newcommand{\eps}{\varepsilon}
\newcommand{\m}{\mathfrak m}
\newcommand{\ot}{\otimes}
\newcommand{\id}{\mathrm{id}}
\newcommand{\unit}{\mathrm e}

\title[A human-scale rank-four counterexample]
      {A Human-Scale Rank-Four Counterexample\\
       from the Affine Group}
\author{}
\date{}
\subjclass[2020]{14L15, 16T05}
\keywords{finite flat group scheme, skew primitive, affine group,
power word, square-zero deformation}

\begin{document}

\begin{abstract}
We give a deliberately mnemonic presentation of a finite locally free
group scheme of rank four which is not killed by four.  Its base is
\[
 R=\Z[a,b]/(a^3,b^3,a^2b+2),
\]
and its coordinate algebra is
\[
 A=R[U,V]/(U^2-abU+b^2V,\ V^2-a^2V).
\]
Put
\[
 \lambda=(1+aU)(1+bV).
\]
The coproduct formulas
\[
 \Delta(U)=U\ot1+\lambda\ot U,\qquad
 \Delta(V)=V\ot\lambda+1\ot V
\]
say geometrically that \((U,V,\lambda)\) defines a closed subgroup of
the standard weight-\((1,-1)\) semidirect product
\(\Ga^2\rtimes\Gm\).  In fact this subgroup assertion has only one
nonformal component.  If the bridge \(b^2V\) is replaced by \(tV\), with
\(t\in(b^2)\), the complete multiplication-defect vector is
\[
 \bigl(a(b^2-t)V_1U_2,0,0\bigr).
\]
Thus \(t=b^2\) kills the only curvature tautologically; all group-law
identities are inherited from the ambient semidirect product.  Modulo the
square-zero ideal \(Q=(b^2)\), the group is the exact matched product of
two rank-two affine subgroups.  Inside the fixed smooth ambient group,
its naive coefficientwise lift has the single normal curvature
\(ab^2V_1U_2\).  The counterexample is obtained by bending just one
defining hypersurface by \(b^2V\); its normal coboundary cancels that
curvature exactly.
Writing \(\lambda=1+\theta\), one finds
\[
 \theta^2=0,\qquad 2\theta=2bV.
\]
Ordinary matrix powering therefore gives
\[
 [4]^\#(U)=2bUV\ne0,\qquad [4]^\#(V)=0,\qquad [8]=\unit.
\]
The nonzero class is transparent from
\[
 R_{\mathrm{add}}\simeq
 (\Z/4)\{1,b\}\oplus
 (\Z/2)\{a,a^2,ab,b^2,ab^2\}.
\]
The same presentation makes the deformation mechanism visible through
the square-zero chain
\[
 b^2\longmapsto ab^2\longmapsto a^2b^2=-2b\ne0.
\]
We first illustrate the mechanism through three elementary deformations:
an affine-line graph, an infinitesimal Heisenberg curvature, and a
characteristic-two finite bridge with one closure defect.
We also explain why an honest product or semidirect product of familiar
smaller group schemes cannot furnish a simpler counterexample.
\end{abstract}

\maketitle

\begin{center}
\fbox{\begin{minipage}{0.92\textwidth}
\small
\textbf{Disclosure: AI-generated work.}\ This document in its entirety---the
construction, the proofs, the computations, and the exposition---was generated
by artificial intelligence. The work was carried out by the AI models Codex
(OpenAI) and Claude (Anthropic).
\end{minipage}}
\end{center}
\medskip

\section{The example}

The analogue of Lagrange's theorem for a finite locally free group scheme
\(G/S\) of rank \(n\) asks whether the power word
\[
 [n]\colon G\longrightarrow G,\qquad g\longmapsto g^n,
\]
is the constant unit word.  The question is recorded in SGA~3; see
\cite[Exp.~VIII, Rem.~7.3.1]{SGA3II}.  It is true over a field and over a
reduced base, and Deligne proved it for commutative finite locally free
group schemes over every base; see \cite[\S1]{OortTate}.  Thus a
counterexample must use both a nonreduced base and a noncommutative group
law.

Here is the construction in full.

\begin{theorem}\label{thm:main}
Let
\begin{equation}\label{eq:R}
 R=\Z[a,b]/(a^3,b^3,a^2b+2)
\end{equation}
and
\begin{equation}\label{eq:A}
 A=R[U,V]/(U^2-abU+b^2V,\ V^2-a^2V).
\end{equation}
Put \(W=UV\) and
\begin{equation}\label{eq:lambda}
 \lambda=(1+aU)(1+bV).
\end{equation}
Then:
\begin{enumerate}
\item \(R\) is a local Artin complete intersection of length nine and
characteristic four, with
\[
 \Soc(R)=\F_2(2b).
\]
\item \(A\) is free of rank four over \(R\), on
\[
 1,\qquad U,\qquad V,\qquad W.
\]
\item There is a Hopf algebra structure on \(A\) with
\(\eps(U)=\eps(V)=0\) and
\begin{equation}\label{eq:Delta}
 \Delta(U)=U\ot1+\lambda\ot U,\qquad
 \Delta(V)=V\ot\lambda+1\ot V.
\end{equation}
\item If \(G=\Spec A\), then
\begin{equation}\label{eq:powers}
 [4]^\#(U)=2bW\ne0,\qquad
 [4]^\#(V)=0,\qquad
 [8]=\unit.
\end{equation}
In particular, the rank-four group scheme \(G\) is not killed by four.
\item Its closed fiber is \(\alpha _2\times\alpha _2\), but \(G\) itself
is noncommutative.
\end{enumerate}
\end{theorem}

The point of this note is not merely that the formulas are short.  We will
derive every term in them from one geometric idea.

The reader interested only in the shortest Hopf-algebra verification can
read Section~\ref{sec:design}, Subsection~\ref{subsec:one-curvature}, and
Section~\ref{sec:power}: ambient multiplication leaves a single curvature
coefficient, which the bridge kills tautologically, and one matrix power
supplies inverses.
The intervening sections explain why that coefficient and that bridge are
forced.

\section{The affine-group design principle}
\label{sec:design}

Let \(M=\Spec B\) be an affine monoid scheme.  A group-like function
\(l\in B\) is exactly a multiplicative character:
\[
 \Delta(l)=l\ot l,\quad\eps(l)=1
 \qquad\Longleftrightarrow\qquad
 l\colon M\longrightarrow\mathbb A^1_{\mathrm{mult}}.
\]
If \(M\) is a group scheme, \(l\) is invertible and this is a homomorphism
\(M\to\Gm\).

Suppose
\[
 \Delta(x)=x\ot1+l\ot x.
\]
On points this says
\[
 x(gh)=x(g)+l(g)x(h).
\]
Thus \(x\) is not usually a homomorphism to \(\Ga\); rather, the pair
\((x,l)\) is a homomorphism to the affine group
\[
 \operatorname{Aff}_1=\Ga\rtimes\Gm,\qquad
 (u,z)(u',z')=(u+zu',zz').
\]
Equivalently,
\[
 g\longmapsto
 \begin{pmatrix}l(g)&x(g)\\0&1\end{pmatrix}
\]
is a two-dimensional representation.  The right skew-primitive formula
\[
 \Delta(y)=y\ot l+1\ot y
\]
has the analogous lower-triangular interpretation.

It is convenient to remember both coordinates simultaneously.  Let
\begin{equation}\label{eq:H-law}
 H=\Ga^2\rtimes\Gm,\qquad
 (u,v,z)(u',v',z')
 =(u+zu',\,vz'+v',\,zz').
\end{equation}
This is the fiber product over \(\Gm\) of an upper and a lower affine
group.  It has the faithful block-matrix representation
\begin{equation}\label{eq:block-matrix}
 (u,v,z)\longmapsto
 \begin{pmatrix}z&u\\0&1\end{pmatrix}
 \oplus
 \begin{pmatrix}z&0\\v&1\end{pmatrix}.
\end{equation}
It is also literally the usual semidirect product for a
\(\Gm\)-action on \(\Ga^2\).  Put
\[
 w=z^{-1}v.
\]
In the coordinates \((u,w,z)\), the law becomes
\begin{equation}\label{eq:standard-semidirect}
 (u,w,z)(u',w',z')
 =
 (u+zu',\,w+z^{-1}w',\,zz').
\end{equation}
Thus
\[
 H\simeq\Ga^2\rtimes\Gm,
 \qquad
 z\cdot(u,w)=(zu,z^{-1}w).
\]
One may represent this standard semidirect product by
\[
 (u,w,z)\longmapsto
 \begin{pmatrix}
  z&0&u\\
  0&z^{-1}&w\\
  0&0&1
 \end{pmatrix}.
\]
The right-skew coordinate \(v\) is simply the denominator-free coordinate
\(v=zw\).

Matrix powering gives the central formula.  If \(l\) is group-like and
\(x\) is left \(l\)-skew primitive, then
\begin{equation}\label{eq:norm}
 [n]^\#(x)=N_n(l)x,\qquad
 N_n(T)=1+T+\cdots+T^{n-1}.
\end{equation}
For a right skew primitive \(y\), one has
\([n]^\#(y)=yN_n(l)\).  Indeed,
\[
 (u,z)^n=(N_n(z)u,z^n)
\]
in \(\operatorname{Aff}_1\).

Now suppose \(4=0\), \(l=1+\theta\), and \(\theta^2=0\).  Then
\begin{equation}\label{eq:carry}
 N_4(l)=2\theta,\qquad l^4=1,\qquad N_8(l)=0.
\end{equation}
Consequently a rank-four counterexample will follow if we can make
\[
 l\text{ group-like},\qquad
 x\text{ left }l\text{-skew primitive},\qquad
 2\theta x\ne0.
\]
The second coordinate in Theorem~\ref{thm:main} is the small amount of
extra memory needed to make a finite rank-four locus in \(H\) stable under
multiplication.

\section{Three toy deformations}
\label{sec:toys}

Before returning to the counterexample, we isolate its three basic ideas
in examples which require almost no machinery.  The first explains why
the affine group is the natural ambient group.  The second explains
curvature and coboundaries.  The third is a small finite flat bridge
construction in which the entire subgroup check has one component.

\subsection{The additive line bending toward the multiplicative line}

Let \(S\) be a ring and let \(\varepsilon\in S\) satisfy
\(\varepsilon^2=0\).  On the affine line define
\begin{equation}\label{eq:toy-affine-law}
 x\star y=x+y+\varepsilon xy.
\end{equation}
This law is best understood without checking associativity directly.
In the affine group
\[
 \operatorname{Aff}_1(S)
 =\{(u,z):z\in S^\times\},
 \qquad
 (u,z)(u',z')=(u+zu',zz'),
\]
consider the graph
\begin{equation}\label{eq:toy-affine-graph}
 j(x)=(x,1+\varepsilon x).
\end{equation}
One calculation gives
\[
 j(x)j(y)
 =
 \bigl(x+y+\varepsilon xy,\,
       (1+\varepsilon x)(1+\varepsilon y)\bigr)
 =
 j(x\star y).
\]
Thus \eqref{eq:toy-affine-law} is associative because it is the
restriction of an already associative multiplication.  Its identity is
\(0\), and
\[
 x^{-1}=-x+\varepsilon x^2.
\]

On coordinate rings this says
\[
 \Delta(x)
 =x\ot1+1\ot x+\varepsilon x\ot x
 =x\ot1+(1+\varepsilon x)\ot x,
\]
while \(1+\varepsilon x\) is group-like.  This is the simplest example
of a skew-primitive coordinate: the apparently deformed addition law is
just a subgroup graph inside an affine group.

There is also a finite version.  Suppose \(S\) has characteristic two.
Then the equation \(x^2=0\) is preserved, since in the tensor square
\[
 \Delta(x)^2
 =x^2\ot1+1\ot x^2+\varepsilon^2x^2\ot x^2=0.
\]
Consequently
\[
 S[x]/(x^2)
\]
with the displayed coproduct is a finite locally free group scheme of
rank two, deforming the usual infinitesimal additive group.  It is still
killed by two because
\[
 [2]^\#(x)=2x+\varepsilon x^2=0.
\]

The general rank-two building block is almost as short.  Impose
\[
 x^2=cx,\qquad \ell=1+dx,
\]
and use \(x_*=x_1+\ell_1x_2\).  Reduction by the two quadratic
relations gives
\begin{equation}\label{eq:toy-rank-two-cancellation}
 x_*^2-cx_*
 =\ell_1(2+cd)x_1x_2,
 \qquad
 1+dx_*=(1+dx_1)(1+dx_2).
\end{equation}
Thus the single scalar equation \(cd=-2\) makes this a rank-two subgroup
of \(\operatorname{Aff}_1\); it also gives
\(\ell^2=1\) and \(x\star x=(2+cd)x=0\).  This is the elementary rank-two cancellation
used twice in the main construction, and is a convenient special case
of the Tate--Oort picture \cite{OortTate}.

There is a useful elementary deformation space hidden here.  Over a
field \(k\) of characteristic two, the parameter equation becomes
\[
 k[c,d]/(cd),
\]
the union of two lines crossing at the origin.  At the origin the group
is
\[
 k[x]/(x^2),\qquad \Delta(x)=x\ot1+1\ot x.
\]
Along the \(c\)-axis the equation of the finite scheme deforms while the
coproduct stays primitive; along the \(d\)-axis the scheme \(x^2=0\)
stays fixed while the coproduct deforms.  Over the dual numbers,
\[
 c=\alpha\varepsilon,\qquad d=\beta\varepsilon
\]
may be chosen independently because \(\varepsilon^2=0\).  Over
\(k[\varepsilon]/(\varepsilon^3)\), however, the first obstruction is
\[
 cd=\alpha\beta\varepsilon^2.
\]
Thus the two directions coexist to first order but obstruct one another
at second order.  The rank-four bridge will provide a new normal
coordinate in which precisely such a mixed obstruction can be canceled.

This toy already contains the group-like character, the skew-primitive
coordinate, and the advantage of inheriting the group axioms from an
ambient semidirect product.

\subsection{A square-zero curvature which may or may not be removable}

Let \(k\) be a nonzero commutative ring and put
\[
 S=k[\varepsilon]/(\varepsilon^2).
\]
On \(\mathbb A^3_S\) define the infinitesimal Heisenberg law
\begin{equation}\label{eq:toy-heisenberg-law}
 (x,y,z)(x',y',z')
 =
 \bigl(x+x',\,y+y',\,z+z'+\varepsilon xy'\bigr).
\end{equation}
Associativity is the single identity
\[
 xy'+(x+x')y''
 =
 x'y''+x(y'+y'').
\]
If \(p=(x,y)\), \(q=(x',y')\), and \(c(p,q)=xy'\), this is the cocycle
identity
\begin{equation}\label{eq:toy-cocycle}
 c(p,q)+c(p+q,r)=c(q,r)+c(p,q+r).
\end{equation}
The inverse is
\[
 (x,y,z)^{-1}=(-x,-y,-z+\varepsilon xy).
\]

Modulo \(\varepsilon\), the plane \(z=0\) is a subgroup.  Its naive lift
is not: multiplying two points of that plane produces the normal error
\[
 \varepsilon xy'.
\]
Try to repair the plane by bending it to
\[
 z=\varepsilon f(x,y).
\]
It is a subgroup exactly when
\begin{equation}\label{eq:toy-coboundary-equation}
 f(x+x',y+y')-f(x,y)-f(x',y')=xy'.
\end{equation}
The left side is symmetric in the two pairs \((x,y)\) and
\((x',y')\), whereas the right side is not.  Hence no such polynomial
\(f\) exists.  The curvature is a cocycle by ambient associativity, but
it is not a coboundary.

For comparison, replace \(xy'\) in
\eqref{eq:toy-heisenberg-law} by the symmetric curvature
\[
 xy'+x'y.
\]
Now
\[
 (x+x')(y+y')-xy-x'y'=xy'+x'y,
\]
so \(f(x,y)=xy\) solves the bending equation.  The graph
\[
 z=\varepsilon xy
\]
is a subgroup, and the coordinate
\[
 \widetilde z=z-\varepsilon xy
\]
turns the entire law back into ordinary addition.  This is the
beginner's deformation-theoretic dictionary:
\[
 \begin{array}{c@{\quad\Longleftrightarrow\quad}l}
  \text{ambient associativity}
    &\text{the curvature is a \(2\)-cocycle},\\
  \text{bending the equations}
    &\text{adding a coboundary},\\
  \text{a subgroup lift exists}
    &\text{the curvature class vanishes}.
 \end{array}
\]

\subsection{A miniature finite bridge}

We next give a small successful bridge which is very close in shape to
the counterexample but has no characteristic-four carry.  Let
\[
 B=\F_2[r,s]/(r^2,s^2),
 \qquad Q=(s).
\]
Thus \(Q^2=0\), and multiplication by \(r\) gives only the two-step chain
\[
 s\longmapsto rs\longmapsto0.
\]
For \(t\in Q\), set
\begin{equation}\label{eq:toy-bridge-algebra}
 A_t
 =
 B[U,V]/(U^2-rU+tV,\ V^2),
 \qquad
 \lambda=(1+rU)(1+sV).
\end{equation}
Use the product coordinates inherited from the denominator-free
semidirect product:
\[
 U_*=U_1+\lambda_1U_2,
 \qquad
 V_*=V_1\lambda_2+V_2.
\]
Reduction by the two quadrics gives the complete defect vector
\begin{equation}\label{eq:toy-bridge-defect}
 \begin{pmatrix}
  U_*^2-rU_*+tV_*\\
  V_*^2\\
  (1+rU_*)(1+sV_*)-\lambda_1\lambda_2
 \end{pmatrix}
 =
 \begin{pmatrix}
  r(s-t)V_1U_2\\0\\0
 \end{pmatrix}.
\end{equation}
Here is a short audit.  The identities \(r^2=s^2=0\) give
\(\lambda_i^2=1\), and hence \(V_*^2=0\).  They also give
\[
 r(1+\lambda_1)=rsV_1,\qquad
 t(1+\lambda_2)=rtU_2,
\]
which reduce the first quadratic defect to the displayed coefficient.
For the graph, put
\[
 \delta=rsV_1U_2,\quad
 A=(1+rU_1)(1+rU_2),\quad
 C=(1+sV_1)(1+sV_2).
\]
Then
\[
 1+rU_*=A+\delta,\qquad 1+sV_*=C+\delta.
\]
Since \(\delta A=\delta C=\delta\) and \(\delta^2=0\), the two errors
cancel in characteristic two:
\[
 (A+\delta)(C+\delta)=AC=\lambda_1\lambda_2.
\]
Thus the naive lift \(t=0\) has the unique curvature
\[
 rsV_1U_2.
\]
There is a concrete group-theoretic interpretation.  Inside the ambient
semidirect product are two rank-two subgroups
\[
 P:\ (U,0,1+rU),\quad U^2=rU,
 \qquad
 K_s:\ (0,V,1+sV),\quad V^2=0.
\]
Their ordered product \(P K_s\) is exactly the scheme defined by
\(A_0\), but a product of two subgroups need not itself be a subgroup.
To reverse the middle factors one would have to send
\[
 U\longmapsto\widetilde U=(1+sV)U.
\]
The failure of \(\widetilde U\) to remain on \(P\) is
\[
 \widetilde U^{\,2}-r\widetilde U=rsVU.
\]
This is precisely the curvature above.  Modulo \((s)\) it vanishes and
the ordered product becomes a matched-product subgroup.  The bridge
will bend this product locus in the normal \(V\)-direction so that it
still closes before reducing modulo \(s\).

The normal cochain \(\eta_t=-tV\) changes it by
\[
 \partial\eta_t
 =-t(V_*-V_1-V_2)
 =-rtV_1U_2.
\]
Choosing \(t=s\) therefore kills the curvature and gives
\begin{equation}\label{eq:toy-successful-bridge}
 A_s=B[U,V]/(U^2-rU+sV,\ V^2).
\end{equation}
This is the exact toy version of the bridge in the main example.

The algebra \(A_s\) is free of rank four on \(1,U,V,UV\).  Moreover
\[
 (1+rU)^2=(1+sV)^2=1,
 \qquad \lambda^2=1.
\]
Hence \(N_4(\lambda)=0\), and the fourth-power word is the identity.
The third-power word supplies inverses, so \(A_s\) is a Hopf algebra
without a separate antipode calculation.  This toy is already a
nontrivial, generally noncommutative deformation, but it is not a
counterexample: it is killed by four.
For example, doubling is visibly nonzero:
\[
 [2]^\#(U)=sUV+rsV,\qquad
 [2]^\#(V)=rUV,
\]
whereas applying doubling once more gives the identity section.

The comparison with the main construction is now visible.  The toy has
only
\[
 s\longmapsto rs\longmapsto0.
\]
The main example replaces this by the longer chain
\[
 q=b^2\longmapsto aq=ab^2
 \longmapsto a^2q=-2b\ne0.
\]
In both cases the middle class is the subgroup curvature and the bridge
is its null-homotopy.  Only in the main example is there one more
nonzero step; that final step is exactly the fourth-power class.

\section{Why an obvious semidirect product cannot work}
\label{sec:no-go}

Before constructing the group, it is useful to understand why the most
obvious interpretation of \eqref{eq:H-law} is too simple.

\begin{lemma}[Killedness propagates through extensions]
\label{lem:extension-no-go}
Suppose
\[
 1\longrightarrow H\longrightarrow G\longrightarrow K\longrightarrow1
\]
is an exact sequence of group schemes.  If \(H\) is killed by \(h\) and
\(K\) is killed by \(k\), then \(G\) is killed by \(hk\).
\end{lemma}

\begin{proof}
For every test algebra and every \(g\in G\), the element \(g^k\) maps to
the identity in \(K\), hence belongs to \(H\).  Therefore
\[
 g^{hk}=(g^k)^h=1.
\]
This pointwise identity is functorial and hence is an identity of word
morphisms.
\end{proof}

In particular, an honest semidirect product of two rank-two group schemes
is killed by four: each rank-two factor is commutative and is killed by
two.  Therefore the desired group cannot retain a flat rank-two kernel
and a flat rank-two quotient.  The bridge term in \eqref{eq:A} is not
cosmetic; it destroys precisely this overly simple extension structure.

Nor does adding another quadratic skew coordinate help the geometric-sum
mechanism.  It doubles the rank from four to eight, while
\eqref{eq:carry} already says \(N_8(l)=0\).  Thus the natural
larger-rank enlargement erases the carry instead of clarifying it.

\section{Two rank-two cancellations and one bridge}
\label{sec:bridge}

We next derive the scalar equations which make a rank-four locus in \(H\)
closed under multiplication.  Begin with
\begin{equation}\label{eq:general-A}
 A_0=B[U,V]/(U^2-\alpha U-\beta V,\ V^2-\gamma V)
\end{equation}
and put
\[
 \lambda=(1+rU)(1+sV).
\]
Try the product law inherited from \(H\):
\begin{equation}\label{eq:general-product}
 U_*=U_1+\lambda_1U_2,\qquad
 V_*=V_1\lambda_2+V_2.
\end{equation}

The elementary rank-two calculation explains the diagonal conditions.
If \(T^2=cT\), \(g=1+dT\), and
\(T_*=T_1+g_1T_2\), then
\begin{equation}\label{eq:rank-two-cancellation}
 T_*^2-cT_*=g_1(2+cd)T_1T_2.
\end{equation}
Thus \(cd=-2\) cancels the cross term.  Applying this to both coordinates,
and then comparing the remaining mixed terms, produces five equations:
\begin{equation}\label{eq:five}
 \begin{array}{c@{\qquad}c@{\qquad}l}
  \alpha r=-2,&\gamma s=-2,
    &\text{two rank-two cancellations},\\[2pt]
  \gamma r=0,&\beta s=0,
    &\text{two mixed vanishings},\\[2pt]
  \multicolumn{2}{c}{\beta r+\alpha s=0}
    &\text{the bridge balance}.
 \end{array}
\end{equation}

\begin{lemma}[The one closure calculation]\label{lem:closure}
Assume \eqref{eq:five}.  In the tensor square of \(A_0\), the product
coordinates \eqref{eq:general-product} satisfy
\begin{equation}\label{eq:three-closure-identities}
 U_*^2=\alpha U_*+\beta V_*,
 \qquad
 V_*^2=\gamma V_*,
 \qquad
 (1+rU_*)(1+sV_*)=\lambda_1\lambda_2.
\end{equation}
\end{lemma}

\begin{proof}
This is a universal reduction by two monic quadrics.  A short hand
calculation is recorded in Appendix~\ref{app:closure}; its only scalar
inputs are the five displayed equations \eqref{eq:five}.
\end{proof}

\begin{proposition}[The bridge criterion]\label{prop:bridge}
Assume \eqref{eq:five}.  Then \(A_0\) is free of rank four over \(B\),
and the formulas
\[
 \Delta(U)=U\ot1+\lambda\ot U,\qquad
 \Delta(V)=V\ot\lambda+1\ot V,\qquad
 \eps(U)=\eps(V)=0
\]
make \(A_0\) a Hopf algebra.  If \(W=UV\), then
\begin{equation}\label{eq:general-power}
 [4]^\#(U)=2sW,\qquad [4]^\#(V)=0,\qquad [8]=\unit.
\end{equation}
\end{proposition}

\begin{proof}
The two monic quadrics give the basis \(1,U,V,W\).  The five scalar
equations imply
\begin{equation}\label{eq:scalar-consequences}
 4=2\alpha=2\beta=2\gamma=2r=0,\qquad
 r^2\beta=2s,\qquad 2s^2=0.
\end{equation}
For example,
\[
 4=(\alpha r)(\gamma s)=\alpha s(\gamma r)=0,
\]
and multiplying the bridge balance by \(r\) gives
\[
 r^2\beta=-r\alpha s=2s.
\]

Let \(L=1+rU\), \(M=1+sV\), and
\(\theta=\lambda-1\).  Reduction by the quadrics gives
\begin{equation}\label{eq:lambda-identities}
 M^2=1,\qquad
 L^2=\lambda^2=1+2sV,\qquad
 \theta^2=0,\qquad
 2\theta=2sV.
\end{equation}
In particular, \(\lambda^{-1}=1-\theta\).  Hence
\[
 j\colon\Spec A_0\longrightarrow H_B,\qquad
 (U,V)\longmapsto(U,V,\lambda)
\]
is a closed immersion: on coordinate rings it sends \(u,v,z\) to
\(U,V,\lambda\), and \(U,V\) generate \(A_0\).  Lemma~\ref{lem:closure}
says exactly that its image is closed under multiplication.  We have
therefore constructed a closed submonoid of the semidirect-product group
\(H\).  Associativity, the unit, and coassociativity of the induced
coproduct are inherited from \(H\); there are no further bialgebra
identities to check.  Inverses will follow from the power calculation.

By \eqref{eq:lambda-identities} and \eqref{eq:carry},
\[
 N_4(\lambda)=2sV.
\]
The norm formula gives
\[
 [4]^\#(U)=2sVU=2sW,
 \qquad
 [4]^\#(V)=2sV^2=2s\gamma V=-4V=0.
\]
It also gives \([8]=\unit\).  Hence the seventh-power word is a two-sided
inverse, so the closed submonoid is a group scheme and \(A_0\) is a Hopf
algebra.
\end{proof}

\section{The three relations are forced}
\label{sec:forced}

We now reverse-engineer the shortest memorable specialization of
Proposition~\ref{prop:bridge}.  Choose the two character parameters to be
\[
 r=a,\qquad s=b.
\]
We want a single scalar relation to supply both diagonal cancellations.
Taking
\[
 a^2b=-2
\]
forces the convenient choices
\[
 \alpha=ab,\qquad \gamma=a^2,
\]
because
\[
 \alpha r=(ab)a=a^2b,\qquad
 \gamma s=(a^2)b=a^2b.
\]
The bridge balance then asks for
\[
 \beta a=-(ab)b,
\]
so the evident choice is
\[
 \beta=-b^2.
\]
The two remaining mixed vanishings become exactly
\[
 \gamma r=a^3=0,\qquad
 \beta s=-b^3=0.
\]

Thus the entire example is forced by
\[
 a^3=0,\qquad b^3=0,\qquad a^2b=-2.
\]
In table form, the five compatibility checks are literally
\begin{equation}\label{eq:specialization-table}
 \begin{array}{c@{\qquad}c}
  \alpha r=(ab)a=a^2b=-2&
  \gamma s=(a^2)b=a^2b=-2\\
  \gamma r=a^3=0&
  \beta s=-b^3=0\\
  \multicolumn{2}{c}{
   \beta r+\alpha s=(-b^2)a+(ab)b=0.}
 \end{array}
\end{equation}
Substituting these coefficients into \eqref{eq:general-A} gives exactly
\eqref{eq:A}.  Proposition~\ref{prop:bridge} now proves all the Hopf
assertions in Theorem~\ref{thm:main}, except for the claimed
nonvanishing.  We prove that next.

In the standard semidirect-product coordinates
\eqref{eq:standard-semidirect}, this closed immersion is
\begin{equation}\label{eq:G-to-standard-semidirect}
 G\longrightarrow\Ga^2\rtimes\Gm,\qquad
 (U,V)\longmapsto
 \bigl(U,\lambda^{-1}V,\lambda\bigr).
\end{equation}
Indeed, if \(Y=\lambda^{-1}V\), then
\[
 \Delta(Y)=Y\ot1+\lambda^{-1}\ot Y.
\]
Thus \(U\) and \(Y\) are the weight \(+1\) and weight \(-1\) translation
coordinates for the \(\Gm\)-action.  Conversely,
\(V=\lambda Y\), so \eqref{eq:G-to-standard-semidirect} is still a closed
immersion and loses no coordinate information.

Equivalently, this gives a construction entirely inside the standard
semidirect product.  Give
\[
 \widetilde H_R=\Ga^2\rtimes\Gm
\]
coordinates \((x,y,z)\), with \(z\) acting with weights \(+1,-1\) as in
\eqref{eq:standard-semidirect}.  Then \(G\) is the closed subscheme cut
out by only three equations:
\begin{equation}\label{eq:three-equations-in-semidirect}
 \begin{split}
  x^2-abx+b^2zy&=0,\\
  zy^2-a^2y&=0,\\
 z-(1+ax)(1+bzy)&=0.
 \end{split}
\end{equation}
Equivalently, for every \(R\)-algebra \(C\), set
\[
 G(C)=\left\{(x,y,z)\in C\times C\times C^\times:
 \begin{array}{l}
  x^2-abx+b^2zy=0,\\
  zy^2-a^2y=0,\\
  z=(1+ax)(1+bzy)
 \end{array}\right\}.
\]
The group law is not an additional formula to discover: it is the
restriction of
\[
 (x,y,z)(x',y',z')
   =(x+zx',\ y+z^{-1}y',\ zz')
\]
from \(\widetilde H(C)\).  Lemma~\ref{lem:closure} is exactly the single
assertion that the three displayed equations are preserved by this
product.  The identity and associativity are then inherited from the
ambient semidirect product, and the already proved identity \([8]=\unit\)
makes the seventh-power word the inverse.  Thus this functor is visibly a
closed subgroup functor of \(\widetilde H_R\).

Finally, \(U=x\), \(V=zy\), and \(\lambda=z\) identify its coordinate
algebra with \(A\); conversely \(y=\lambda^{-1}V\).  This is the
semidirect-product construction of the group scheme.  Notice that it
appears to require three closure identities, rather than separate
verifications of the coassociativity, counit, and antipode axioms.  The
next subsection shows that two of those identities are automatic and
the third has only one scalar coefficient.

Here ``semidirect-product construction'' means a closed subgroup of an
ambient semidirect product; it does not mean that \(G\) itself splits as
a semidirect product of finite group schemes.  Lemma~\ref{lem:extension-no-go}
shows that such a splitting could not give a counterexample.

\subsection{The one-curvature form of the construction}
\label{subsec:one-curvature}

The semidirect-product presentation reduces the entire construction to
one closure coefficient.  Put
\[
 q=b^2,\qquad Q=(q)\subset R.
\]
Then
\[
 Q^2=bQ=2Q=0.
\]
For any \(t\in Q\), consider
\begin{equation}\label{eq:At}
 A_t=
 R[U,V]/(U^2-abU+tV,\ V^2-a^2V),
 \qquad
 \lambda=(1+aU)(1+bV).
\end{equation}
The cases \(t=0\) and \(t=q\) will be the naive lift and the corrected
lift, respectively.

\begin{lemma}[The one-curvature calculation]\label{lem:one-curvature}
In \(A_t\ot_R A_t\), put
\[
 \lambda_i=(1+aU_i)(1+bV_i),\qquad
 U_*=U_1+\lambda_1U_2,\qquad
 V_*=V_1\lambda_2+V_2.
\]
Then the complete closure-defect vector is
\begin{equation}\label{eq:one-curvature-vector}
 \begin{pmatrix}
  U_*^2-abU_*+tV_*\\
  V_*^2-a^2V_*\\
  (1+aU_*)(1+bV_*)-\lambda_1\lambda_2
 \end{pmatrix}
 =
 \begin{pmatrix}
  a(q-t)V_1U_2\\0\\0
 \end{pmatrix}.
\end{equation}
\end{lemma}

\begin{proof}
This is one reduction by the two monic quadrics; a short hand proof is
given in Appendix~\ref{app:one-curvature}.  Conceptually, the two diagonal
rank-two cancellations are both the relation \(a^2b=-2\), the two pure
mixed terms die because \(a^3=0\) and \(bt=0\), and only the coefficient
\[
 a(q-t)
\]
can remain.  The important extra content of
\eqref{eq:one-curvature-vector} is that the other two closure equations
vanish identically, before imposing this last coefficient.
\end{proof}

Now take \(t=q=b^2\).  The right side of
\eqref{eq:one-curvature-vector} vanishes tautologically.  Thus the map
\[
 \Spec A_q\longrightarrow H_R,
 \qquad
 (U,V)\longmapsto(U,V,\lambda)
\]
identifies its source with a closed submonoid of the ambient semidirect
product.  The identities in Appendix~\ref{app:one-curvature} also show
that \(\lambda\) is a unit.  Associativity, the identity, coassociativity,
the counit, and group-likeness of \(\lambda\) are inherited from \(H\).
Finally, \([8]=\unit\) makes the seventh-power word the inverse, so this
closed submonoid is a group scheme.  Since \(A_q=A\), this gives the
shortest proof of the Hopf assertion in Theorem~\ref{thm:main}: apart from
the power calculation needed for inverses, the only nonformal
verification is the single coefficient \(a(b^2-b^2)=0\).

In the standard weight-\((1,-1)\) coordinates \((x,y,z)\), the same family
is cut out by
\[
 x^2-abx+tzy=0,\qquad zy^2-a^2y=0,\qquad
 z=(1+ax)(1+bzy).
\]
Under ambient multiplication its sole defect is
\[
 a(q-t)z_1y_1x_2.
\]

\section{The base ring and the surviving top class}
\label{sec:base}

The base presentation has a transparent additive calculation.  Start with
\[
 R_0=\Z[a,b]/(a^3,b^3),
\]
which is free over \(\Z\) on the nine box monomials
\[
 a^ib^j,\qquad 0\le i,j\le2.
\]
Let \(N\) be multiplication by \(a^2b\).  It has only two nonzero arrows:
\begin{equation}\label{eq:two-arrows}
 1\longmapsto a^2b,\qquad
 b\longmapsto a^2b^2.
\end{equation}
The additive group of \(R\) is the cokernel of \(N+2\id\).
Each arrow in \eqref{eq:two-arrows} gives a \(\Z/4\)-summand, while each
of the other five box monomials gives a \(\Z/2\)-summand.  Consequently
\begin{equation}\label{eq:additive-R}
 R_{\mathrm{add}}\simeq
 (\Z/4)\{1,b\}\oplus
 (\Z/2)\{a,a^2,ab,b^2,ab^2\}.
\end{equation}
Equivalently, every element has a unique expression
\[
 c_0+c_1b+
 \epsilon_1a+\epsilon_2a^2+\epsilon_3ab+
 \epsilon_4b^2+\epsilon_5ab^2,
\]
where \(c_0,c_1\in\Z/4\) and \(\epsilon_i\in\F_2\).

This proves at once that \(R\) has length nine and that
\begin{equation}\label{eq:nonzero-carry}
 2b=-a^2b^2\ne0.
\end{equation}
It also shows \(4=0\).  Alternatively,
\[
 4=(a^2b)^2=a^4b^2=0.
\]
Since \(a,b\) are nilpotent and \(2=-a^2b\), the ring is local with
maximal ideal \(\m=(a,b)\) and residue field \(\F_2\).  Its associated
graded ring is
\[
 \gr_{\m}R\simeq\F_2[A,B]/(A^3,B^3),
\]
so its Hilbert function is \((1,2,3,2,1)\).  The top class is
\[
 A^2B^2\longleftrightarrow a^2b^2=-2b,
\]
and hence
\[
 \Soc(R)=\F_2(2b).
\]
Finally, \(a^3,b^3,a^2b+2\) form a regular sequence in
\(\Z[a,b]\), for example by the injectivity of \(N+2\id\) in the
calculation above.  Thus \(R\) is a complete intersection.

The two relations defining \(A\) are successively monic, so
\[
 A=R\{1,U,V,W\}.
\]
Combining this freeness with \eqref{eq:nonzero-carry} proves
\[
 2bW\ne0.
\]

\section{The fourth power is visible in one matrix}
\label{sec:power}

Put
\[
 \theta=\lambda-1=aU+bV+abW.
\]
The bridge criterion gives
\begin{equation}\label{eq:theta}
 \theta^2=0,\qquad 2\theta=2bV.
\end{equation}
Therefore
\[
 N_4(\lambda)=2\theta=2bV.
\]
The universal point of \(G\), viewed through
\eqref{eq:block-matrix}, has fourth power
\begin{equation}\label{eq:matrix-fourth}
 \left[
 \begin{pmatrix}\lambda&U\\0&1\end{pmatrix}
 \oplus
 \begin{pmatrix}\lambda&0\\V&1\end{pmatrix}
 \right]^4
 =
 \begin{pmatrix}1&2bUV\\0&1\end{pmatrix}
 \oplus I_2.
\end{equation}
The upper-right entry is nonzero by
\eqref{eq:nonzero-carry} and freeness of \(A\).  This single matrix is
the counterexample:
\[
 \operatorname{rank}(G)=4,\qquad [4]\ne\unit.
\]
On the other hand, \(N_8(\lambda)=0\), so the eighth power of the same
matrix is the identity.

Modulo \(\m=(a,b)\), both quadrics become \(U^2=V^2=0\), \(\lambda=1\),
and both coordinates are primitive.  Thus
\[
 G_{\F_2}\simeq\alpha _2\times\alpha _2.
\]
The failure is entirely relative.  The total coproduct is nevertheless
noncocommutative: in
\[
 \Delta(U)-\tau\Delta(U)
 =(\lambda-1)\ot U-U\ot(\lambda-1),
\]
the coefficient of \(V\ot U\) is \(b\ne0\).

\section{The counterexample as a bent subscheme}
\label{sec:square-zero}

We now recast the construction as an embedded first-order deformation.
This is perhaps its most geometric form.  A simple subgroup downstairs
is cut out by three hypersurfaces in the fixed smooth group \(H\).  The
coefficientwise lift of those hypersurfaces is no longer a subgroup: the
product of two points misses one hypersurface by one square-zero normal
error.  Bending that hypersurface by \(b^2V\) cancels the error exactly.

Put
\[
 q=b^2,\qquad Q=(q)\subset R,\qquad D=R/Q.
\]
Then
\[
 Q^2=0,\qquad bQ=2Q=0,
\]
and multiplication by \(a\) gives the chain
\begin{equation}\label{eq:memory-chain}
 q=b^2\longmapsto aq=ab^2
 \longmapsto a^2q=a^2b^2=-2b
 \longmapsto0.
\end{equation}

The closed immersion
\[
 \Spec D\longrightarrow\Spec R
\]
is therefore a first-order thickening.  It is important that it is not a
split dual-number thickening; we return to this point below.

\subsection{The simple subscheme downstairs}

Use the denominator-free coordinates of \eqref{eq:H-law} and write
\[
 B_R=\mathcal O(H_R)=R[U,V,z,z^{-1}],\qquad
 \lambda=(1+aU)(1+bV).
\]
Consider the three ambient equations
\begin{equation}\label{eq:three-normal-equations}
 \Phi=U^2-abU,\qquad
 \Psi=V^2-a^2V,\qquad
 \Xi=z-\lambda.
\end{equation}
After reduction to \(D\), let
\[
 G_0=V(\overline\Phi,\overline\Psi,\overline\Xi)
 \subset H_D.
\]
This is the simpler object from which the counterexample will be lifted.

Modulo \(Q\), the bridge disappears.  There are two rank-two subgroup
schemes of \(H_D\):
\begin{equation}\label{eq:two-downstairs-factors}
 \begin{array}{lll}
  P:&U^2=abU,&p(U)=(U,0,1+aU),\\[2pt]
  K:&V^2=a^2V,&k(V)=(0,V,1+bV).
 \end{array}
\end{equation}
Both subgroup assertions are the same elementary rank-two cancellation,
because
\[
 (ab)a=(a^2)b=a^2b=-2.
\]
Multiplication in \(H_D\) gives
\begin{equation}\label{eq:matched-product-map}
 P\times_DK\longrightarrow H_D,
 \qquad
 (U,V)\longmapsto p(U)k(V)
   =\bigl(U,V,(1+aU)(1+bV)\bigr).
\end{equation}
Its image is \(G_0=G_D\), and the map is an isomorphism of schemes onto
that image.  It is closed under multiplication by
Lemma~\ref{lem:one-curvature} with \(t=0\), since \(q=0\) over \(D\).
Moreover \(2b=-a^2q=0\) over \(D\), so the norm calculation gives
\([4]=\unit\); hence it is a subgroup.  Thus \(G_D\) is an \emph{exact
matched product} \(P\bowtie K\).  It need not be a direct or semidirect
product as a group; reversing the two factors produces mutual actions.
All coherence for those actions is inherited from associativity in \(H\).

There is a useful geometric bonus: \(G_0\subset H_D\) is a complete
intersection with completely explicit normal directions.  Indeed,
\(\overline\Phi\) and \(\overline\Psi\) are successive monic equations.
After imposing them, \(\overline\lambda^2=1\), so
\(z-\overline\lambda\) is a nonzerodivisor in the remaining Laurent
polynomial ring.  If \(I_0\) denotes the ideal of \(G_0\), then
\begin{equation}\label{eq:conormal-basis}
 I_0/I_0^2\simeq
 \mathcal O(G_0)\{\overline\Phi,\overline\Psi,\overline\Xi\}.
\end{equation}
Thus an infinitesimal embedded displacement may be recorded simply by
saying how far each of these three equations moves.

\subsection{The naive lift and its one normal curvature}

Lift the three equations in \eqref{eq:three-normal-equations}
coefficientwise from \(D\) to \(R\):
\begin{equation}\label{eq:naive-embedded-lift}
 X_{\mathrm{naive}}=V(\Phi,\Psi,\Xi)\subset H_R.
\end{equation}
Eliminating \(z\) and reducing by the two monic quadrics shows that this
is a finite flat rank-four lift of \(G_0\), with basis
\(1,U,V,UV\).  It contains the ambient identity, but it is not closed
under ambient multiplication.

Here is an intrinsic way to measure that failure.  Put
\[
 E=Q\ot_D\mathcal O(G_0),\qquad
 E_2=Q\ot_D\mathcal O(G_0\times_DG_0).
\]
The equality \(Q^2=0\), together with flatness, identifies \(E\) with
the square-zero ideal in the coordinate ring of any one of our lifts.
Since \(G_0\) is a subgroup, pulling a lifted defining equation back
under ambient multiplication and then restricting it to
\(X_{\mathrm{naive}}\times_RX_{\mathrm{naive}}\) gives an element of
\(E_2\).  The result depends only on its conormal class and defines the
\emph{multiplication curvature}
\[
 \kappa_{\mathrm{naive}}:
 I_0/I_0^2\longrightarrow E_2,
\]
where the target is viewed as an \(\mathcal O(G_0)\)-module through the
multiplication of \(G_0\).  Lemma~\ref{lem:one-curvature} says that the
whole map is
\begin{equation}\label{eq:naive-curvature}
 \begin{split}
  \kappa_{\mathrm{naive}}(\overline\Phi)
    &=aq\,V_1U_2,\\
  \kappa_{\mathrm{naive}}(\overline\Psi)
    &=\kappa_{\mathrm{naive}}(\overline\Xi)=0.
 \end{split}
\end{equation}
In plain language, the product of two points remains on the second
quadratic hypersurface and on the character graph, but misses the first
quadratic hypersurface by the infinitesimal normal distance
\(aqV_1U_2\).  This is the only failure of the naive lift.

\subsection{Bending one hypersurface}

Because \(Q^2=0\), correcting the first equation is a linear problem.
For \(t\in Q\), define
\begin{equation}\label{eq:embedded-lift-family}
 X_t=V(\Phi+tV,\Psi,\Xi)\subset H_R.
\end{equation}
Every \(X_t\) is finite flat of rank four and has special fiber \(G_0\).
Fixing \(X_{\mathrm{naive}}\) as an origin, its equation displacement is
the conormal homomorphism
\begin{equation}\label{eq:normal-equation-displacement}
 \begin{split}
  \delta_t:I_0/I_0^2&\longrightarrow E,\\
  \overline\Phi&\longmapsto t\otimes\overline V,\qquad
  \overline\Psi,\overline\Xi\longmapsto0.
 \end{split}
\end{equation}
This is the standard normal-module description of the difference
between two framed embedded lifts across a square-zero extension.  It is
an affine difference, not a canonical vector attached to a lift: the
naive lift has supplied the origin.

There is a harmless but useful sign convention.  The equation is
\[
 \Phi+tV=0,
 \qquad\text{or geometrically}\qquad
 \Phi=-tV.
\]
Thus the equation displacement is \(+tV\), whereas the geometric normal
displacement is the one-cochain
\[
 \eta_t=-tV
\]
in the \(\overline\Phi\)-direction.  The linearized subgroup condition
applied to this geometric displacement is especially simple.  Using
\(V_*=V_1\lambda_2+V_2\), one obtains the single-line calculation
\begin{equation}\label{eq:bridge-coboundary}
 \begin{split}
  d_{\mathrm{mult}}\eta_t
  &=-t(V_*-V_1-V_2)\\
  &=-tV_1(\lambda_2-1)
   =-atV_1U_2.
 \end{split}
\end{equation}
Here \(bt=0\) removes the \(V_2\)-part of
\(t(\lambda_2-1)\).  Hence the subgroup-lifting equation is simply
\[
 \kappa_{\mathrm{naive}}+d_{\mathrm{mult}}\eta_t
 =a(q-t)V_1U_2=0.
\]
This is the active component of the general linearized
multiplication-stability operator on the normal module.  We do not need
its other components: \eqref{eq:naive-curvature} says that they vanish.
Because \(Q^2=0\), this first-order calculation is exact; there are no
discarded quadratic or higher-order deformation terms.

The most evident solution is \(t=q=b^2\), and it gives exactly
\[
 U^2-abU+b^2V=0.
\]
Consequently
\begin{equation}\label{eq:counterexample-as-bend}
 G=X_q=V(\Phi+qV,\Psi,\Xi)\subset H_R
\end{equation}
is obtained from the matched-product subgroup \(G_0\) by one first-order
normal bend inside the fixed smooth ambient group \(H_R\).  Equivalently,
in the standard weight-\((1,-1)\) coordinates \(V=zy\), we have changed
only the first wall
\[
 x^2-abx=0
 \qquad\text{to}\qquad
 x^2-abx=-b^2zy.
\]

The bend makes \(X_q\) a closed submonoid and does not move the unit.
In fact no separate antipode formula is needed here.

\begin{lemma}[Inverses lift in a finite flat monoid]
\label{lem:monoid-lifts-group}
Let \(S_0\hookrightarrow S\) be a nilpotent closed immersion and let
\(M/S\) be a finite locally free monoid scheme.  If \(M_{S_0}\) is a
group scheme, then \(M\) is a group scheme.
\end{lemma}

\begin{proof}
Consider the left-translation morphism over the first factor
\[
 M\times_SM\longrightarrow M\times_SM,
 \qquad (x,y)\longmapsto(x,xy).
\]
It is a morphism between finite locally free schemes of the same rank
over \(M\).  Its reduction over \(S_0\) is an isomorphism.  Locally on
\(M\), the determinant of the corresponding map of finite free modules
is therefore a unit modulo a nilpotent ideal, hence a unit; the original
morphism is an isomorphism.  The same argument applies to right
translation.  Pulling the unit section back under the two inverse
translation maps supplies right and left inverses.  Associativity makes
them equal.
\end{proof}

Apply the lemma to \(G=X_q\) and \(G_0=G_D\).  Thus multiplicative
closure, finite flatness, and the simple group downstairs already imply
that the bent subscheme is a group scheme.  The independently computed
identity \([8]=\unit\) gives the more explicit inverse word \([7]\).

Ambient associativity also forces the usual cocycle identity for the
curvature: multiplying three points in the two possible parenthesizations
gives the same next obstruction.  Thus the bridge may be viewed as a
null-homotopy of the unique multiplication curvature.  The displayed
two-term calculation is the relevant one-coordinate piece of the usual
bialgebra deformation complex of Gerstenhaber--Schack
\cite{GerstenhaberSchack}; it is not ordinary extension cohomology for
the two rank-two factors.

\subsection{The subgroup locus and the fourth-power locus}

One may package the preceding construction in the embedded deformation
functor of finite flat degree-four subschemes of \(H\), or, informally,
in the relative Hilbert functor \(\operatorname{Hilb}^4(H)\).  The point
\([G_0]\) lies over \(D\).  Requiring the universal subscheme to contain
the unit and to be stable under multiplication and inversion cuts out
the subgroup subfunctor.  Requiring in addition that its fourth-power
word equal the unit cuts out a smaller subfunctor.  Our calculation says
that the naive lift misses the first locus by \(aqV_1U_2\), while the
normal bend \(-qV\) moves it back into that locus.  We now see that the
resulting group still leaves the second locus.

Indeed, let \(A=\mathcal O(G)\), put \(I=QA\), and write
\(A_D=A/I=\mathcal O(G_0)\).  Freeness of \(A\) gives
\[
 I^2=0,\qquad I\simeq Q\ot_DA_D.
\]
Let \(c:G\to G\) be the constant unit map.  Since
\(2b=-a^2q=0\) over \(D\), the power calculation gives
\([4]_{G_0}=c_{G_0}\).  Hence the difference of the two algebra maps
\([4]^\#-c^\#\) lands in \(I\), kills \(I\), and induces a map
\[
 \omega_4:A_D\longrightarrow Q\ot_DA_D.
\]
The difference of two algebra maps which agree modulo a square-zero ideal
is a derivation relative to their common reduction.  Thus \(\omega_4\)
is a \(c_{G_0}^\#\)-derivation, and the exact power calculation gives
\begin{equation}\label{eq:first-order-word-defect}
 \omega_4(\overline U)
   =2b\,\overline U\overline V
   =-a^2q\,\overline U\overline V\ne0,
 \qquad
 \omega_4(\overline V)=0.
\end{equation}
Equivalently, it is the first-order displacement of the fourth-power map
from the unit map,
\[
 c_{G_0}^*\Omega^1_{G_0/D}
 \longrightarrow Q\ot_D\mathcal O_{G_0}.
\]
It should not be confused with the differential of \([4]\) in the group
variable; it is the variation of the identity \([4]=c\) in the base
deformation.  It is nonzero because \(2b\ne0\) in \(Q\) and
\(1,U,V,UV\) is a free basis upstairs.

The whole mechanism can now be read as a three-line geometric story:
\begin{equation}\label{eq:three-stage-deformation-story}
 \underbrace{qV}_{\text{normal bend}}
 \quad\rightsquigarrow\quad
 \underbrace{aq\,V_1U_2}_{\substack{\text{multiplication curvature}\\
                                      \text{canceled by the bend}}}
 \quad\rightsquigarrow\quad
 \underbrace{-a^2q\,UV=2bUV}_{\substack{\text{fourth-word defect}\\
                                         \text{which survives}}}.
\end{equation}
More precisely, in this chosen normal coordinate,
\[
 \omega_4(\overline U)
 =-a\,\operatorname{diag}^*
   \bigl(\kappa_{\mathrm{naive}}(\overline\Phi)\bigr).
\]
This is the coefficient-memory chain \eqref{eq:memory-chain}.  A
square-zero coefficient module need not be one-dimensional or have
trivial scalar action: \(Q\) remembers \(q,aq,a^2q\) in succession.
That is why a genuinely first-order deformation can cancel the subgroup
obstruction and still retain a secondary word obstruction.

The possible corrections in this normal slice form the torsor
\begin{equation}\label{eq:bridge-torsor}
 \{t\in Q:a(q-t)=0\}
 =q+\operatorname{Ann}_Q(a)
 =\{q,\ q+2b\}.
\end{equation}
Indeed, the chain \eqref{eq:memory-chain} gives
\(\operatorname{Ann}_Q(a)=(a^2q)=(2b)\).  Both choices have the same
secondary residue
\[
 -a^2t=2b.
\]
The norm calculation turns this residue into
\([4]^\#(U)=2bUV\ne0\).  In other words, the primary subgroup curvature
is a coboundary, but its null-homotopy leaves a nonzero power-word class
one step farther along the \(a\)-chain.

The extension \(R\to D\) is visibly nonsplit.  Every lift \(b+x\), with
\(x\in Q\), of \(\overline b\in D\) satisfies
\[
 (b+x)^2=b^2\ne0,
\]
whereas \(\overline b^{\,2}=0\) in \(D\).  The construction is therefore
a first-order lift across a general nonsplit square-zero thickening, not
literally a dual-number tangent arc.  The nonsplitting is also the source
of the curvature: the coefficient \(\overline b\) squares to zero
downstairs, while every lift has the unavoidable square \(q\) upstairs.
Thus the inhomogeneous bridge is not an artifact of a poor choice of
coordinates.

There is also a direct framed gauge check.  A coordinate change which is
the identity modulo \(Q\) has the form
\[
 U'=U+q\phi,\qquad V'=V+q\psi.
\]
Since \(q^2=2q=bq=0\), one has
\[
 (U')^2-abU'+qV'
 =U^2-abU+qV+q\phi(2U-ab)
 =U^2-abU+qV.
\]
Thus the bridge cannot be removed by an infinitesimal coordinate change
within this natural normal slice.

\section{Why this seems close to the simplest human model}

There are two different notions of simplicity.

The fully universal bridge base is
\[
 B_{\mathrm{br}}
 =
 \frac{\Z[\alpha,\beta,\gamma,r,s]}
 {(\alpha r+2,\gamma s+2,\gamma r,\beta s,\beta r+\alpha s)}.
\]
It is not Artinian, and every defining equation has exactly one geometric
job.  It is the best presentation for discovering the construction.
The map
\[
 B_{\mathrm{br}}\longrightarrow R,\qquad
 (\alpha,\beta,\gamma,r,s)
 \longmapsto(ab,-b^2,a^2,a,b)
\]
then gives the two-variable model above.  It is the best presentation for
remembering a concrete counterexample.

For comparison, the optimized paper uses
\[
 \Z[r,s]/(r^3+2,s^3+2,rs^2)
\]
and the coefficients \(r^2,-rs,s^2\).  The present base is not even
larger numerically: both bases have length nine and use three relations.
The gain is explanatory.  Here both diagonal cancellations are the same
equation \(a^2b=-2\), both mixed vanishings are pure cubic nilpotence, and
the bridge balance is a two-term monomial syzygy.

The two base directions are not gratuitous.  In any local target of the
bridge equations with \(2s\ne0\), neither \(\alpha\) nor \(r\) can be a
unit, while \(\alpha r=-2\).  Hence
\[
 2\in\m^2.
\]
This already rules out the tempting bases
\((\Z/4)[t]/(t^N)\), for which \(2\notin(2,t)^2\).  It also rules out
local Artin chain rings more generally.  Indeed, let \(v\) be the chain
valuation, let \(N\) be the nilpotence length, and put \(e=v(2)\).
The equations give
\[
 v(\alpha)+v(r)=e,\qquad
 v(\gamma)+v(s)=e,\qquad
 v(\gamma)+v(r)\ge N,\qquad
 v(\beta)+v(s)\ge N.
\]
The bridge is nonzero when \(2s\ne0\), and hence
\[
 v(\beta)+v(r)
 =v(\alpha)+v(s)=e-v(r)+v(s).
\]
The third inequality gives
\[
 v(r)-v(s)\ge N-e>0,
\]
whereas the fourth gives
\[
 e+2\bigl(v(s)-v(r)\bigr)
 =v(\beta)+v(s)\ge N,
\]
a contradiction.  Two transverse nilpotent directions are therefore
structurally natural in this design.

More group coordinates do not remove the essential difficulty.  Honest
extensions are excluded by Lemma~\ref{lem:extension-no-go}, while the
first obvious extra skew coordinate changes the rank to eight and makes
the norm vanish.  The pair \((U,V)\) is therefore not an optimization
artifact: \(U\) detects the affine-group carry and \(V\) is the minimal
memory coordinate which lets a finite rank-four locus close under
multiplication.

The example can be remembered in four lines:
\[
 \boxed{
 \begin{gathered}
 a^3=b^3=0,\qquad a^2b=-2,\\
 U^2=abU-b^2V,\qquad V^2=a^2V,\\
 \lambda=(1+aU)(1+bV),\\
 [4]^\#(U)=2bUV\ne0,\qquad [8]=\unit.
 \end{gathered}}
\]

\appendix

\section{The one-curvature calculation}\label{app:one-curvature}

We prove Lemma~\ref{lem:one-curvature}.  Write
\[
 L=1+aU,\qquad M=1+bV,\qquad \lambda=LM.
\]
The base relations and \(t\in Q=(b^2)\) imply
\[
 2a=4=0,\qquad t^2=bt=2t=0,
 \qquad a^2b^2=-2b.
\]
Reduction by \(U^2=abU-tV\) and \(V^2=a^2V\) gives
\begin{equation}\label{eq:one-curvature-auxiliary}
 \begin{gathered}
  M^2=1,\qquad
  \lambda^2=1-a^2tV,\qquad
  a^2(\lambda-1)=-2V,\\
  t(\lambda-1)=atU,\qquad
  t(\lambda^2-1)=0,\\
  2\lambda U+ab(\lambda^2-\lambda)=-ab^2V.
 \end{gathered}
\end{equation}
For the last identity, the two summands reduce respectively to
\[
 2U+2bUV
 \quad\text{and}\quad
 2U-ab^2V+2bUV;
\]
their sum is \(-ab^2V\) because \(4=0\).

The second quadratic now has no defect:
\[
 V_*^2-a^2V_*
 =V_1\lambda_2
   \bigl(a^2(\lambda_2-1)+2V_2\bigr)=0.
\]
For the first quadratic, direct expansion followed by
\eqref{eq:one-curvature-auxiliary} gives
\begin{align*}
 U_*^2-abU_*+tV_*
 &=tV_1(\lambda_2-1)\\
 &\quad+
 \bigl(2\lambda_1U_1+ab(\lambda_1^2-\lambda_1)\bigr)U_2
 +t(1-\lambda_1^2)V_2\\
 &=a(t-b^2)V_1U_2
  =a(q-t)V_1U_2.
\end{align*}
The last equality uses \(2a=0\).

It remains to check the graph coordinate.  This calculation has a useful
factor-swapping form.  Put
\[
 A=aU_2,\qquad B=bV_1,\qquad E=AB.
\]
Then
\[
 1+aU_*=L_1(1+A+E),
 \qquad
 1+bV_*=M_2(1+B+E),
\]
and hence
\begin{align*}
 &(1+aU_*)(1+bV_*)-\lambda_1\lambda_2\\
 &\hspace{35mm}
 =L_1M_2E(2+A+B+E).
\end{align*}
Now \(2E=0\).  Moreover
\[
 B^2=-2B,
 \qquad
 A^2=-a^2tV_2,
\]
so \(EB=EA=E^2=0\), using \(bt=0\).  The graph defect is therefore
zero.  This completes the exact three-component calculation
\eqref{eq:one-curvature-vector}.

\section{The universal closure calculation}\label{app:closure}

We record the calculation behind Lemma~\ref{lem:closure}.  It is the only
nonformal step in passing from the ambient semidirect-product group to the
finite rank-four subgroup.

The five equations \eqref{eq:five} imply
\[
 4=2\alpha=2\beta=2\gamma=2r=0,\qquad
 r^2\beta=2s,\qquad 2s^2=0.
\]
They also give
\[
 \gamma(\lambda-1)=-2V,\qquad
 \beta(1-\lambda)=\alpha sU,\qquad
 \beta(\lambda^2-1)=0
\]
and
\[
 2\lambda U+\alpha(\lambda^2-\lambda)+\alpha sV=0.
\]
Consequently
\begin{align*}
 V_*^2-\gamma V_*
 &=V_1\lambda_2
   \bigl(\gamma(\lambda_2-1)+2V_2\bigr)=0,\\
 U_*^2-\alpha U_*-\beta V_*
 &=\beta V_1(1-\lambda_2)\\
 &\quad+
 \bigl(2\lambda_1U_1+\alpha(\lambda_1^2-\lambda_1)\bigr)U_2
 +\beta(\lambda_1^2-1)V_2=0.
\end{align*}
This proves preservation of the two quadrics.

For the graph coordinate, let
\[
 L=1+rU,\qquad M=1+sV,\qquad E=rsV_1U_2.
\]
Direct factorization gives
\[
 (1+rU_*)(1+sV_*)-\lambda_1\lambda_2
 =L_1M_2E(2+rU_2+sV_1+E).
\]
Each of the four terms on the right vanishes by the scalar consequences
above and the two quadrics.  This proves
\[
 (1+rU_*)(1+sV_*)=\lambda_1\lambda_2
\]
and finishes the closure calculation.

\section{Exact computer audit}

All arguments above are by hand.  A short independent Macaulay2 script
checks the polynomial identities by exact strong Gr\"obner reduction over
\(\Z\):
\begin{center}
 \path{m2/verify_human_scale_affine_rank4_counterexample_20260712.m2}.
\end{center}
From the project root, run
\begin{verbatim}
M2 --script m2/verify_human_scale_affine_rank4_counterexample_20260712.m2
\end{verbatim}
The script checks \(4=0\), \(2b\ne0\), and the nonzero normal form of
\(2bUV\).  It verifies that the coproduct preserves both quadrics, that
\(\lambda\) is group-like, that the coproduct is noncocommutative, and
that the displayed antipode satisfies both convolution identities.  In
universal coefficients, it verifies that the five bridge equations imply
all three identities of Lemma~\ref{lem:closure}.  In the square-zero
normal slice, it independently obtains the exact curvature vector
\[
 \bigl(a(t-b^2)V_1U_2,0,0\bigr)
\]
and verifies that the nonzero naive curvature is canceled by the
coboundary \(-b^2(V_*-V_1-V_2)\).  It separately checks that
\[
 q,\ aq,\ a^2q\ne0,\qquad q^2=bq=2q=0,
\]
that the naive character satisfies \(\lambda^2=1\), and that the
first-order fourth-word defect obeys the conceptual identity
\[
 \omega_4(U)
 =-a\,\operatorname{diag}^*
   \bigl(\kappa_{\mathrm{naive}}(\Phi)\bigr).
\]
In the standard weight-\((1,-1)\) semidirect coordinates, it independently
checks that all three equations
\eqref{eq:three-equations-in-semidirect} are closed under multiplication.
It also checks
\[
 \theta^2=0,\qquad 2\theta=2bV,\qquad
 ([4]^\#U,[4]^\#V,[8]^\#U,[8]^\#V)=(2bUV,0,0,0),
\]
as well as the naive square-zero curvature before the bridge is inserted.
Finally, it audits the three beginner examples: the affine-line graph
and its rank-two quotient, associativity of the infinitesimal Heisenberg
cocycle, and the complete defect vector
\[
 \bigl(r(s-t)V_1U_2,0,0\bigr)
\]
for the miniature bridge.  In the corrected toy it checks
\(\lambda^2=1\), the two displayed nonzero doubling formulas,
\(N_4(\lambda)=0\), and noncocommutativity.

\begin{thebibliography}{9}

\bibitem{SGA3II}
M.~Demazure and A.~Grothendieck (eds.),
\emph{Sch\'emas en groupes II: Groupes de type multiplicatif, et structure
des sch\'emas en groupes g\'en\'eraux (SGA 3)},
Lecture Notes in Mathematics, vol.~152, Springer-Verlag, 1970.

\bibitem{GerstenhaberSchack}
M.~Gerstenhaber and S.~D.~Schack,
\emph{Bialgebra cohomology, deformations, and quantum groups},
Proc. Natl. Acad. Sci. USA \textbf{87} (1990), no.~1, 478--481,
\href{https://doi.org/10.1073/pnas.87.1.478}
{doi:10.1073/pnas.87.1.478}.

\bibitem{OortTate}
J.~Tate and F.~Oort,
\emph{Group schemes of prime order},
Ann. Sci. \'Ecole Norm. Sup. (4) \textbf{3} (1970), no.~1, 1--21.

\end{thebibliography}

\end{document}
